\scanpage{26}
\[
\Delta(N_{\Omega/S})=-\deleted{\unsure{4}}\unsure{\mathscr D|\mathcal D|\mathfrak D}(\Omega/S) \tag*{(19)}
\]
Pour ceci, on va d'abord calculer la forme bil.
$\varphi$ associée \unsure{à|a} $N$, on trouve
\[
\varphi_{\Omega/S}(\zeta,\eta)=\operatorname{tr}\zeta\operatorname{tr}\eta-\operatorname{tr}(\zeta\eta) \tag*{(20)}
\]
[en particulier
\[
\varphi_{\Omega/S}(\zeta,1)=\operatorname{tr}\zeta\qquad ] \tag*{(21)}
\]
i.e.
\[
N(\zeta+\eta)=N(\zeta)+N(\eta)+(\operatorname{tr}\zeta\operatorname{tr}\eta-\operatorname{tr}(\zeta\eta)).\tag*{(22)}
\]
Pour vérifier la formule (21), on peut choisir un
splitting et calculer avec $b,c$. On peut aussi
la considérer comme \inserted{cas particulier} d'une formule plus
générale, valable pour $\zeta,\eta\in\operatorname{End}(B)$, où $B$
est un fibré vectoriel de rang 2 quelconque\ldots

On conclut de (20) \deleted{que} la forme discriminant
\deleted{est} donnée par
\[
{\textstyle\bigwedge^2}\varphi(\zeta\wedge\zeta',\eta\wedge\eta')=\det
\begin{pmatrix}\varphi(\zeta,\eta)&\varphi(\zeta,\eta')\\
\varphi(\deleted{\unsure{\eta}}\zeta',\eta)&\varphi(\deleted{\unsure{\eta}}\zeta',\eta')\end{pmatrix}
\]
on va prendre $\zeta'=e_B$, $\eta'=e_B$, \inserted{$\zeta=\eta$} et on trouve
\[
\begin{aligned}
{\textstyle\bigwedge^2}\varphi( e_B\wedge\zeta,e_B\wedge\zeta)
 &=\det\begin{pmatrix}(\operatorname{tr}\zeta)^2-\operatorname{tr}(\zeta^2)&\operatorname{tr}\zeta\\
 \operatorname{tr}\zeta&2\deleted{\unsure{\zeta}}\end{pmatrix}
 =\deleted{\unsure{2}}(\operatorname{tr}\zeta)^2-2\operatorname{tr}(\zeta^2)\\
 &\deleted{=(\operatorname{tr}\zeta)^2\operatorname{tr}\unsure{\eta}-(\operatorname{tr}\eta)^2}
\end{aligned}
\]
On veut vérifier que \unsure{$-$} c'est égal à $-\deleted{\unsure{4}}\mathscr D\cdot u^2\deleted{v}$, où
$u,v$ sont les images de $\zeta,\deleted{\eta}$ dans $\Gamma\check L$. Pour ceci,
\deleted{\unsure{obs|OBS|OPS}} que l'on a un scindage $s$, et que
$\zeta$ \deleted{\unsure{y est de}} \inserted{est} de la forme $su$, \deleted{\unsure{on}} et on a
\deleted{trouve,} comme $\zeta^2=s(u)^2=-cu^2\oplus(bu^2)$, $\operatorname{tr}\zeta^2=-2cu^2+b^2u^2$
\deleted{trouve}
\[
\begin{aligned}
{\textstyle\bigwedge^2}\varphi(1\wedge s(u),1\wedge s(u))&=\deleted{\unsure{(bu)^2-bu}}(bu)^2-2\deleted{\unsure{\operatorname{tr}\zeta^2}}(-2cu^2+b^2u^2)\\
&=(4c-b^2)u^2=-\mathscr D u^2,\quad \text{cqfd.}
\end{aligned}
\]

% archive sheet number: 25

